% !TeX root = ../sections/02_exercises.tex
\begin{exercise}
	数列 $\{a_n\}$ と $\{b_n\}$ はそれぞれ
	$\alpha\in\mathbb{R}$ および
	$\beta\in\mathbb{R}\setminus\{0\}$ に収束するとする。
	以下の問いに答えよ。
	
	\begin{enumerate}
		\item $N_0\in\mathbb{N}$ が存在して，
		$n\geq N_0$ ならば
		\[
		|b_n|\geq \frac{|\beta|}{2}
		\]
		であることを証明せよ。
		
		\item 数列
		\[
		\left\{\frac{a_n}{b_n}\right\}_{n\geq N_0}
		\]
		は収束して，
		\[
		\lim_{n\to\infty}\frac{a_n}{b_n}
		=
		\frac{\alpha}{\beta}
		\]
		であることを証明せよ。
	\end{enumerate}
\end{exercise}

\begin{background}
	この問題では，商の極限を扱う。
	分母にあたる数列 $b_n$ が $0$ に近づいてしまうと
	商を扱いにくいので，まず $b_n$ が十分大きな番号以降で
	$0$ から離れていることを示す。
	
	\begin{definition}[数列の収束]
		数列 $\{b_n\}$ が実数 $\beta$ に収束するとは，
		任意の $\epsilon>0$ に対して，ある自然数 $N$ が存在して，
		すべての $n\geq N$ に対して
		\[
		|b_n-\beta|<\epsilon
		\]
		が成り立つことをいう。
	\end{definition}
	
	\begin{remark}
		今回，$\beta\neq 0$ である。
		したがって
		\[
		|\beta|>0
		\]
		である。
		
		$b_n\to\beta$ なので，十分大きな $n$ では
		$b_n$ は $\beta$ に近い。
		そのため，$b_n$ も $0$ から離れているはずである。
	\end{remark}
	
	\subsection{分母を $0$ から離すための準備}
	
	\begin{lemma}[逆三角不等式]
		任意の実数 $x,y$ に対して
		\[
		\bigl||x|-|y|\bigr|\leq |x-y|
		\]
		が成り立つ。
	\end{lemma}
	
	\begin{remark}
		逆三角不等式を $x=b_n,\ y=\beta$ に適用すると，
		\[
		\bigl||b_n|-|\beta|\bigr|
		\leq |b_n-\beta|
		\]
		となる。
		
		したがって，$|b_n-\beta|$ を十分小さくすれば，
		$|b_n|$ は $|\beta|$ に近くなる。
	\end{remark}
	
	\begin{lemma}[商の誤差評価]
		$\beta\neq 0$ かつ $|b_n|\geq |\beta|/2$ のとき，
		\[
		\left|
		\frac{a_n}{b_n}
		-
		\frac{\alpha}{\beta}
		\right|
		\leq
		\frac{2}{|\beta|}|a_n-\alpha|
		+
		\frac{2|\alpha|}{|\beta|^2}|b_n-\beta|
		\]
		が成り立つ。
	\end{lemma}
	
	\begin{remark}
		この評価により，商の誤差を
		\[
		|a_n-\alpha|
		\quad\text{と}\quad
		|b_n-\beta|
		\]
		の誤差に分けて考えることができる。
	\end{remark}
\end{background}
\begin{strategy}
	まず，(1) では $b_n$ が十分大きな番号以降で $0$ から離れていることを示す。
	
	$\beta\neq 0$ であるから
	\[
	|\beta|>0
	\]
	である。
	
	$b_n\to\beta$ なので，極限の定義において
	\[
	\epsilon=\frac{|\beta|}{2}
	\]
	と選ぶことができる。
	
	すると，ある自然数 $N_0$ が存在して，すべての $n\geq N_0$ に対して
	\[
	|b_n-\beta|<\frac{|\beta|}{2}
	\]
	が成り立つ。
	
	逆三角不等式より
	\[
	\bigl||b_n|-|\beta|\bigr|
	\leq |b_n-\beta|
	\]
	であるから，$|b_n|$ は $|\beta|$ から大きく離れない。
	したがって
	\[
	|b_n|>\frac{|\beta|}{2}
	\]
	が得られる。特に
	\[
	|b_n|\geq\frac{|\beta|}{2}
	\]
	である。
	
	次に，(2) では
	\[
	\frac{a_n}{b_n}\to\frac{\alpha}{\beta}
	\]
	を示す。
	
	示すべきことは，任意の $\epsilon>0$ に対して，
	十分大きな $n$ で
	\[
	\left|
	\frac{a_n}{b_n}
	-
	\frac{\alpha}{\beta}
	\right|<\epsilon
	\]
	となることである。
	
	差を通分すると，
	\[
	\frac{a_n}{b_n}-\frac{\alpha}{\beta}
	=
	\frac{\beta a_n-\alpha b_n}{\beta b_n}
	\]
	である。
	
	分子を
	\[
	\beta a_n-\alpha b_n
	=
	\beta(a_n-\alpha)+\alpha(\beta-b_n)
	\]
	と分解する。
	
	すると三角不等式より，
	\[
	|\beta a_n-\alpha b_n|
	\leq
	|\beta||a_n-\alpha|
	+
	|\alpha||b_n-\beta|
	\]
	となる。
	
	また，(1) より十分大きな $n$ では
	\[
	|b_n|\geq \frac{|\beta|}{2}
	\]
	であるから，分母は $0$ から離れている。
	
	したがって，商全体の誤差は
	\[
	|a_n-\alpha|
	\quad\text{と}\quad
	|b_n-\beta|
	\]
	を十分小さくすることで小さくできる。
	
	つまり，(2) の方針は，
	まず (1) によって分母 $b_n$ を $0$ から離し，
	そのうえで商の差を通分し，
	収束の仮定
	\[
	a_n\to\alpha,
	\qquad
	b_n\to\beta
	\]
	を使って誤差を $\epsilon$ より小さく抑えることである。
\end{strategy}